% Generated draft: requires mathematical and source review.
\documentclass[11pt]{article}
\usepackage[margin=25mm]{geometry}
\usepackage{fontspec}
\setmainfont{Latin Modern Roman}
\usepackage{amsmath,amssymb,mathtools,unicode-math}
\setmathfont{Latin Modern Math}
\usepackage{enumitem,needspace,xurl,hyperref}
\setlength{\parindent}{0pt}
\setlength{\parskip}{5pt}
\setlength{\emergencystretch}{3em}
\newcommand{\N}{\mathbb{N}}
\newcommand{\Z}{\mathbb{Z}}
\newcommand{\Q}{\mathbb{Q}}
\newcommand{\R}{\mathbb{R}}
\newcommand{\C}{\mathbb{C}}
\newcommand{\F}{\mathbb{F}}
\newcommand{\A}{\mathbb{A}}
\newcommand{\PP}{\mathbb{P}}
\newcommand{\E}{\mathbb{E}}
\newcommand{\eps}{\varepsilon}
\newcommand{\dd}{\,\mathrm{d}}
\begin{document}
\section*{UnsolvedMath ID: 30000756}
\textbf{Source problem code:} catalogue entry
\subsection{Definitions and mathematical statement}
Let $p$ be any prime and let $(a_n)_{n\ge1}$ be any sequence of digits with
\[
 a_n\in\{0,1,\ldots,p-1\}.
\]
Using exactly the same digits and starting index, define
\[
 \alpha=\sum_{n=1}^{\infty}a_np^{-n}\in\R,
 \qquad
 \alpha_p=\sum_{n=1}^{\infty}a_np^n\in\Q_p.
\]
The first sum converges in the real absolute value and the second in the $p$-adic absolute value. In particular $\alpha\in[0,1]$ and $\alpha_p\in p\Z_p$. The positive exponent in the second sum is $n$, not $n-1$.

The conjecture is
\[
 \bigl(\alpha\text{ algebraic over }\Q\ \text{and}\
       \alpha_p\text{ algebraic over }\Q\bigr)
 \quad\Longleftrightarrow\quad
 \bigl(\alpha\in\Q\ \text{and}\ \alpha_p\in\Q\bigr)
\]
for every prime and every digit sequence above. In the $p$-adic clause, $\Q$ is its natural embedded copy in $\Q_p$. Algebraicity means that each number is a root of a nonzero polynomial in $\Q[x]$; the two numbers need not satisfy the same polynomial or be equal to one another.

The conjecture concerns simultaneous algebraicity, not a claim that either number is individually forced to be rational whenever it alone is algebraic. The digit sequence is fixed as data, so no canonical-choice convention for a nonunique real expansion is required. Eventually zero or eventually $p-1$ tails and the endpoint real value one are not excluded. No symmetry, complexity, recurrence, or periodic-pattern hypothesis on the digits is added.
\subsection{Short English statement}
For the real and $p$-adic numbers formed from the same digits, must simultaneous algebraicity force both numbers to be rational? The claim applies to every prime and every digit sequence, with no pattern restriction.
\subsection{Sources}
\begin{itemize}
\item Local UnsolvedMath registry, permanent ID 30000756, code catalogue entry, ``Simultaneous Real and $p$-Adic Algebraicity.'' Proposed year 2007; proposer unspecified.
\item \emph{Diophantische Approximationen} (2008), Oberwolfach Reports, pp.\ 1115--1190, DOI 10.4171/owr/2007/21; \url{https://doi.org/10.4171/owr/2007/21}. The displayed report year 2008 differs from the DOI and proposed year 2007.
\item Supplied ``Adamczewski--Bugeaud, symmetric real and p-adic expansions'': \url{https://doi.org/10.1093/imrn/rnl111}.
\item Statement status corrected\_verified; the supplied note says the entire equivalence was preserved while digit set, sequence, sums, and exponents were restored.
\item Literature triage checked 2026-08-21: open. It says the simultaneous problem is explicitly unresolved in the report and later symmetric-pattern results do not establish the general criterion. Exact partial-case hypotheses and current status require independent verification. Created 2026-08-25; updated 2026-09-16.
\item Registry curation by Ulam AI / UnsolvedMath Contributors, CC BY 4.0; underlying report and literature provenance remain as supplied. Unreviewed statement-only draft.
\end{itemize}
\end{document}
